%!TeX root=../tese.tex
%("dica" para o editor de texto: este arquivo é parte de um documento maior)
% para saber mais: https://tex.stackexchange.com/q/78101

% As palavras-chave são obrigatórias, em português e em inglês, e devem ser
% definidas antes do resumo/abstract. Acrescente quantas forem necessárias.
\palavraschave{reality gap, benchmark, metric, divergence, robot}

\keywords{Keyword1,Keyword2,Keyword3}

% O resumo é obrigatório, em português e inglês. Estes comandos também
% geram automaticamente a referência para o próprio documento, conforme
% as normas sugeridas da USP.
\resumo{
O campo da manipulação robótica enfrenta a falta de benchmarks acessíveis
 e padronizados para avaliar a transição de modelos do ambiente virtual
  para o físico. Embora o treinamento simulado seja mais seguro,
   rápido e econômico, ele sofre com o reality gap, 
   a divergência técnica e dinâmica entre a simulação e o mundo real. 
   O principal desafio abordado é a escassez de métricas e 
   frameworks capazes de quantificar e mitigar essa disparidade 
   sem exigir hardwares caros. Com isso, busca-se entender como 
   mapear essas inconsistências virtuais para garantir que os 
   robôs desempenhem suas tarefas com precisão no mundo real.
}

\abstract{
Elemento obrigatório, elaborado com as mesmas características do resumo em
língua portuguesa. De acordo com o Regimento da Pós-Graduação da USP (Artigo
99), deve ser redigido em inglês para fins de divulgação. É uma boa ideia usar
o sítio \url{www.grammarly.com} na preparação de textos em inglês.
Text text text text text text text text text text text text text text text text
text text text text text text text text text text text text text text text text
text text text text text text text text text text text text text text text text
text text text text text text text text text text text text.
Text text text text text text text text text text text text text text text text
text text text text text text text text text text text text text text text text
text text text.
}
